\chapter{Dizionari e tavole}
\renewcommand{\arraystretch}{1.5}
\begin{xltabular}{\textwidth}{|>{\raggedright\arraybackslash}p{3cm}|>{\raggedright\arraybackslash}X|>{\raggedright\arraybackslash}X|>{\raggedright\arraybackslash}p{3.5cm}|}
    \caption{Dizionario delle Entità}
    \label{tab:dizionario} \\
    \hline
    \textbf{Entità} & \textbf{Descrizione} & \textbf{Attributi} & \textbf{Identificatore} \\
    \hline
    Pagamento & Pagamenti delle lezioni & ID\_pagamento, \newline Data\_Pagamento, Importo & \textbf{ID\_Pagamento} \\
    \hline
    Studente & Allievo della scuola di musica & Nome, Cognome, Codice\_fiscale, Livello & \textbf{Codice\_Fiscale} \\
    \hline
    Prodotto & Prodotto presente in magazzino & Quantità, \newline Categoria\_\allowbreak Merceologica, Cod\_Prodotto, Nome & \textbf{Cod\_Prodotto} \\
    \hline
    Strumento & Strumento musicale personale o della scuola & ID\_Strumento, Noleggiato & \textbf{ID\_Strumento} \\
    \hline
    Categoria\_\allowbreak Strumento & Categoria dello strumento (e.g a fiato, a corde etc.) & Categoria\_\allowbreak Strumento & \textbf{Categoria\_\allowbreak Strumento} \\
    \hline
    Docente & Insegnante che svolge lezioni di musica & ID\_Docente, Nome, Cognome, Specializzazione & \textbf{ID\_Docente} \\
    \hline
    Workshop & Lezione o lezioni di gruppo con ospiti & Data, ID\_Workshop & \textbf{ID\_Workshop} \\
    \hline
    Aula & Luogo dove si svolgono le lezioni & Strumento\_fisso, ID\_Aula, Numero\_Posti & \textbf{ID\_Aula} \\
    \hline
    Lezione & Attività didattica con Docente & Data, ID\_Lezione & \textbf{ID\_Lezione} \\
    \hline
    Stipendio & Importo mensile dei docenti & Importo, \newline Data\_Mensilità, Docente, Cod\_Pagamento & \textbf{Cod\_Pagamento} \\
    \hline
\end{xltabular}
\newpage
\renewcommand{\arraystretch}{1.5}
\begin{xltabular}{\textwidth}{|>{\raggedright\arraybackslash}p{3cm}|>{\raggedright\arraybackslash}X|>{\raggedright\arraybackslash}X|}
\caption{Dizionario delle Relazioni}
    \label{tab:dizionario} \\
    \hline
    \textbf{Relazione} & \textbf{Partecipanti} & \textbf{Descrizione} \\
    \hline
    Salda & Studente, Pagamento & Lo studente salda il pagamento \\
    \hline
    Noleggio & Studente, Strumento & Associa gli strumenti noleggiati agli studenti \\
    \hline
    Acquista & Studente, Prodotto & Lo studente acquista un prodotto dal magazzino \\
    \hline
    Appartiene & Strumento, Categoria Strumento & Associa ad ogni strumento la sua categoria \\
    \hline
    Si suona & Strumento, Lezione & Uno strumento si suona a lezione \\
    \hline
    Occupa & Lezione, Aula & Associa ogni aula a una lezione
    \hline
    Si svolgono & Aula, Workshop & I workshop si svolgono in un'aula \\
    \hline
    Partecipa & Workshop, docente & Un docente partecipa al workshop \\
    \hline
    Partecipa & Studente, Lezione & Uno studente partecipa alla lezione \\ 
    \hline
    Insegna & Docente, Lezione & Un Docente svolge una lezione
    \hline
\end{xltabular}
\newpage
\renewcommand{\arraystretch}{1.5}
\begin{xltabular}{\textwidth}{|>{\raggedright\arraybackslash}p{3cm}|>{\raggedright\arraybackslash}X|>{\raggedright\arraybackslash}X|}
\caption{Tavola delle frequenze}
    \label{tab:dizionario} \\
    \hline
    \textbf{Operazione} & \textbf{Descrizione} & \textbf{Frequenza} \\
    \hline
    OP1 & Registrazione allievo & 2 volte/mese \\
    \hline
    OP2 & Calcolo stipendi & 15 volte/mse \\ 
    \hline
    OP3 & Registrazione nuova lezione & 100 volte/mese \\ 
    \hline
    OP4 & Acquisto di accessori & 50 volte/mese \\ 
    \hline
    OP5 & Aggiornamento prodotti magazzino & 70 volte/mese \\ 
    \hline
    OP6 & Registrazione di nuovi workshop & 10 volte/mese \\ 
    \hline
    OP7 & Prenotazione aule & 100 volte/mese \\ 
    \hline
    OP8 & Noleggio strumenti & 10 volte/mese \\ 
    \hline
    OP9 & Controllo pagamenti lezioni & 300 volte/mese \\ 
    \hline
    OP10 & verifica disponibilità aula & 400 volte/mese \\ 
    \hline
\end{xltabular}
\newpage
\renewcommand{\arraystretch}{1.5}
\begin{xltabular}{\textwidth}{|>{\raggedright\arraybackslash}p{4cm}|c|c|}
\caption{Tavola dei volumi, stima dopo un anno}
    \label{tab:dizionario} \\
    \hline
    \textbf{Concetto} & \textbf{Tipo} & \textbf{Volume} \\
    \hline
    Pagamento & E & 2000 \\
    \hline
    Studente & E & 200 \\
    \hline
    Prodotto & E & 20 \\ 
    \hline
    Strumento & E & 350 \\ 
    \hline
    Categoria Strumento & E & 5 \\ 
    \hline
    Docente & E & 20 \\ 
    \hline
    Workshop & E & 50 \\ 
    \hline
    Aula & E & 10 \\ 
    \hline
    Lezione & E & 9246 \\ % 200 studenti 1 lezione a settimana (46 settimane) e 1 lezione di gruppo da 8 persone
    \hline
    Stipendio & E & 220 \\ 
    \hline
    Salda & R & 1460 \\ 
    \hline
    Noleggio & R & 30 \\ 
    \hline
    Acquista(Studente, Prodotto) & R & 300 \\ 
    \hline
    Si suona & R & 9246 \\ 
    \hline
    Occupa \newline (Lezione, Aula) & R & 9246 \\ 
    \hline
    Partecipa \newline (Docente, Workshop) & R & 40 \\ 
    \hline
    Partecipa \newline (Studente, Lezione) & R & 9246 \\ 
    \hline
    Insegna \newline (Docente, Lezione) & R & 9286 \\ 
    \hline
\end{xltabular}
\renewcommand{\arraystretch}{1.5}
\begin{xltabular}{\textwidth}{|C|>{\raggedright\arraybackslash}p{3cm}|>{\raggedright\arraybackslash}p{3cm}|}
\caption{Tavola delle frequenze}
 \label{tab:dizionario} \\
    \hline
    \textbf{Operazione} & \textbf{Descrizione} & \textbf{Frequenza} \\
    \hline
    OP1 & Registrazione Allievo & 5/mese \\
    \hline
    OP2 & Calcolo stipendi & 20/mese \\ 
    \hline
    OP3 & Registrazione nuova lezione & 600/mese \\ 
    \hline
    OP4 & Acquisto di accessori & 200/mese \\ 
    \hline
    OP5 & Aggiornamento prodotti magazzino & 100/mese \\ 
    \hline
    OP6 & Registrazione nuovi workshop & 10/mese \\ 
    \hline
    OP7 & Prenotazione aule & 600/mese \\
    \hline
    OP8 & Noleggio strumenti & 20/mese \\ 
    \hline
    OP9 & Controllo pagamenti & 1200/mese \\ 
    \hline
    OP10 & Verifica disponibilità & 1500/mese \\ 
    \hline
\end{xltabular}
\clearpage
\begin{figure}[p]
    \centering
    % NOTA: angle=90 deve stare PRIMA di width e height!
    \includegraphics[angle=0, width=\textwidth, height=0.9\textheight, keepaspectratio]{Schemi ER progetto-Schema OP3.jpg}
    \caption{Schema operazione OP3}
    \label{fig:schema_er}
\end{figure}
\begin{figure}[p]
    \centering
    % NOTA: angle=90 deve stare PRIMA di width e height!
    \includegraphics[angle=0, width=\textwidth, height=0.9\textheight, keepaspectratio]{Schemi ER progetto-Schema OP6.jpg}
    \caption{Schema operazione OP6}
    \label{fig:schema_er}
\end{figure}
\begin{figure}[p]
    \centering
    % NOTA: angle=90 deve stare PRIMA di width e height!
    \includegraphics[angle=0, width=\textwidth, height=0.9\textheight, keepaspectratio]{Schemi ER progetto-Schema OP8.jpg}
    \caption{Schema operazione OP8}
    \label{fig:schema_er}
\end{figure}